\section{Perturbing a Hamiltonian involution}\label{sec:involution}

We construct Hamiltonian maps whose fixed points are prescribed by Morse
functions on the fixed components of an involution.  The involution used in
our example acts on a K\"ahler surface and fixes four projective lines.  Its
product with the identity on another manifold will therefore replace one
Morse problem on the product by four Morse problems on smaller manifolds.

We use the basic Hamiltonian constructions recalled in \cite{McDuffSalamon2017}. Throughout this section the Hamiltonian vector field of a function $f$ on a
symplectic manifold $(P,\omega)$ is defined by
\[
  \omega(X_f,\,\cdot\,)=df.
\]
We first record the elementary period estimate that makes the perturbation
count exact.  This is a short-period estimate of the same kind as
Yorke's Lipschitz bound~\cite{Yorke1969}; the weaker estimate below suffices.

\begin{lemma}\label{lem:short-period}
Let $V$ be a smooth vector field on a closed smooth manifold $P$, and let
$\psi^s$ denote its flow.  There is a number $\tau>0$ such that, for every
$0<T<\tau$,
\[
  \psi^T(x)=x \quad\Longrightarrow\quad V(x)=0.
\]
\end{lemma}

\begin{proof}
Choose finitely many coordinate balls with compact closures inside their
charts, together with smaller balls whose closures lie in the corresponding
coordinate balls and whose union covers $P$.  Bounds on the coordinate
speeds give a time $\delta>0$ such that a
trajectory starting in any smaller ball stays in its larger ball for time
$\delta$.  In these finitely many balls choose a common positive Lipschitz
bound $L$ for the coordinate representatives of $V$.

Suppose a trajectory of period $T<\delta$ starts in one of the smaller balls.
Write its coordinate representation as $y:[0,T]\to\R^d$ and the coordinate
vector field as $v$.  Set
\[
  m=\max_{0\le t\le T}|v(y(t))|,
\]
and choose $t_0$ at which this maximum is attained.  Since $y(T)=y(0)$,
\[
  \int_0^T v(y(t))\,dt=0.
\]
Moreover, $|y(t)-y(t_0)|\le Tm$ for every $t$.  It follows that
\[
\begin{split}
  m
  &=\left|\frac1T\int_0^T
      \bigl(v(y(t_0))-v(y(t))\bigr)\,dt\right|\\
  &\le \frac LT\int_0^T|y(t_0)-y(t)|\,dt
   \le LTm.
\end{split}
\]
Taking $\tau=\min\{\delta,1/L\}$ forces $m=0$ whenever $0<T<\tau$.
\end{proof}

\begin{proposition}\label{prop:involution}
Let $(P,\omega)$ be a closed symplectic manifold, and let $g:P\to P$ be a
Hamiltonian diffeomorphism satisfying $g^2=\mathrm{id}_P$.  Write
\[
  \Fix(g)=F_1\sqcup\cdots\sqcup F_r
\]
for its connected fixed components.  For each $j$, choose a smooth Morse
function $h_j:F_j\to\R$.  There exists a $g$-invariant smooth function
$f:P\to\R$ with $f|_{F_j}=h_j$ such that, for all sufficiently small
$\epsilon>0$, the map
\[
  \Phi_\epsilon=\psi^\epsilon\circ g,
\]
where $\psi^s$ is the Hamiltonian flow of $f$, satisfies
\begin{equation}\label{eq:involution-fixed-set}
  \Fix(\Phi_\epsilon)
    =\bigsqcup_{j=1}^r\operatorname{Crit}(h_j).
\end{equation}
Every fixed point of $\Phi_\epsilon$ is nondegenerate.  Moreover,
$\Phi_\epsilon$ is the time-one map of a smooth one-periodic Hamiltonian.
If $P$ is simply connected, all its fixed-point loops for this Hamiltonian
are contractible.
\end{proposition}

\begin{proof}
We first describe the geometry of the fixed locus and extend the component
functions.  We then identify every fixed point of the perturbation and
check its derivative.

Average a Riemannian metric over the group $\{\mathrm{id}_P,g\}$.  At a fixed
point $x$, its exponential chart intertwines $g$ with $d_xg$.  Thus the fixed
locus is a smooth closed submanifold, with tangent space
\[
  E_x^+=\ker(d_xg-I).
\]
It has finitely many connected components by compactness and this local
description.  Since $g$ is symplectic, the eigenspace decomposition
\[
  T_xP=E_x^+\oplus E_x^-,
  \qquad E_x^-=\ker(d_xg+I),
\]
is symplectically orthogonal: for $u\in E_x^+$ and $v\in E_x^-$,
\[
  \omega_x(u,v)
    =\omega_x(d_xg(u),d_xg(v))=-\omega_x(u,v).
\]
Nondegeneracy of $\omega_x$ then implies that its restrictions to both
summands are nondegenerate.  Each $F_j$ is therefore symplectic, and
$E_x^-$ is its symplectic normal space.

Extend the functions $h_j$ to a smooth function $\widetilde f$ on $P$, using
disjoint tubular neighborhoods of the closed components.  Define
\[
  f=\tfrac12(\widetilde f+\widetilde f\circ g).
\]
Then $f\circ g=f$ and $f|_{F_j}=h_j$.  Invariance and the symplectic property
of $g$ imply $g_*X_f=X_f$, so $g$ commutes with $\psi^s$.  At a point
$x\in F_j$, invariance also gives $d_xf(v)=0$ for $v\in E_x^-$.  Consequently
$X_f$ is tangent to $F_j$, and its restriction is the Hamiltonian vector
field of $h_j$ for $\omega|_{F_j}$.  In particular,
\begin{equation}\label{eq:ambient-equilibria}
  x\in F_j,\quad X_f(x)=0
  \quad\Longleftrightarrow\quad d_xh_j=0.
\end{equation}
Indeed, at a critical point of $h_j$, the differential of $f$ vanishes on
both summands $E_x^+$ and $E_x^-$.

Let $\tau$ be supplied by Lemma~\ref{lem:short-period} for $X_f$.  Commutation
and $g^2=\mathrm{id}_P$ give
\[
  \Phi_\epsilon^2=\psi^{2\epsilon}.
\]
If $2\epsilon<\tau$ and $\Phi_\epsilon(x)=x$, then
$\psi^{2\epsilon}(x)=x$, so the lemma implies $X_f(x)=0$.  Hence
$\psi^\epsilon(x)=x$, and the original fixed-point equation gives $g(x)=x$.
Equation~\eqref{eq:ambient-equilibria} now shows that $x$ is one of the
prescribed component critical points.  Every such point is conversely
fixed by both maps.  This proves~\eqref{eq:involution-fixed-set}.

It remains to check full nondegeneracy.  At a point in this finite fixed
set, let $A_x:T_xP\to T_xP$ be the linearization of $X_f$.  The point is an
equilibrium, so
\[
  d_x\psi^\epsilon=\exp(\epsilon A_x).
\]
Equivariance implies that $A_x$ commutes with $d_xg$ and preserves $E_x^+$
and $E_x^-$.  Write its restrictions as $A_x^+$ and $A_x^-$.  On
$E_x^+=T_xF_j$, the map $A_x^+$ is the linearization of the Hamiltonian
vector field of $h_j$.  The nonsingular Hessian of $h_j$ and the
nondegenerate form $\omega|_{F_j}$ make $A_x^+$ invertible.  Therefore
\[
  d_x\Phi_\epsilon|_{E_x^+}=\exp(\epsilon A_x^+)
\]
has no eigenvalue $1$ for all sufficiently small positive $\epsilon$:
its eigenvalues are $e^{\epsilon\lambda}$ with $\lambda\ne0$.  On the normal
space,
\[
  d_x\Phi_\epsilon|_{E_x^-}=-\exp(\epsilon A_x^-)
      \longrightarrow -I
      \qquad\text{as }\epsilon\longrightarrow0,
\]
so this block also has no eigenvalue $1$ for small $\epsilon$.  One choice
of $\epsilon$ works for the finite set of component critical points.

Finally, choose a smooth Hamiltonian $K_s$, $0\le s\le1$, whose time-one
map is $g$.  Choose nonnegative smooth functions $a,b$ supported in
successive disjoint subintervals of $(0,1)$, each with integral $1$, and
put
\[
  A(t)=\int_0^t a(u)\,du,
  \qquad
  H_\epsilon(t,x)=a(t)K_{A(t)}(x)+\epsilon b(t)f(x).
\]
During the first interval, the flow traverses the path of $K_s$; during
the second, it traverses $\psi^s$ from $s=0$ to $s=\epsilon$.  Its time-one
map is thus $\psi^\epsilon\circ g$.  The function $H_\epsilon$ vanishes near
$t=0$ and $t=1$, so its periodic extension is smooth.  If $P$ is simply
connected, every fixed-point loop along this specified path is
contractible.
\end{proof}

We now construct a surface to which the proposition applies.  The crucial
feature is that its half-turn fixes entire exceptional curves, including
those on which the full circle action is nontrivial.

\begin{proposition}\label{prop:surface}
Let $Y$ be the complex blow-up of $\CP^1\times\CP^1$ at the four points
\[
  \{0,\infty\}\times\{0,\infty\}.
\]
There is a circle-invariant K\"ahler form $\omega_Y$ for which the lifted
diagonal circle action is Hamiltonian.  Its half-turn $g_Y$ has precisely
four fixed components, the exceptional curves $E_1,\ldots,E_4$, each
isomorphic to $\CP^1$.  Along each $E_j$, the derivative of $g_Y$ is the
identity on the tangent bundle and minus the identity on the symplectic
normal bundle.  The surface $Y$ is simply connected, and
\begin{equation}\label{eq:Y-homology}
  H_i(Y;\Z)\cong
  \begin{cases}
    \Z,&i=0,4,\\
    \Z^6,&i=2,\\
    0,&\text{otherwise}.
  \end{cases}
\end{equation}
\end{proposition}

\begin{proof}
On the standard affine chart $\C^2\subset\CP^1\times\CP^1$, let the circle act by
\[
  t\cdot(z,w)=(e^{2\pi it}z,e^{2\pi it}w),
  \qquad t\in\R/\Z.
\]
The four centers are invariant, so the action lifts holomorphically to
the blow-up.  The surface $Y$ is K\"ahler and simply connected by
Proposition~\ref{prop:blowup}.  Averaging a K\"ahler form over the
holomorphic circle action gives an invariant K\"ahler form $\omega_Y$.
If $V$ denotes the infinitesimal generator of the circle action, invariance
gives
\[
  d(\omega_Y(V,\,\cdot\,))=\mathcal L_V\omega_Y=0.
\]
Since $H^1(Y;\R)=0$, this closed one-form is $dk$ for a smooth function
$k:Y\to\R$.  Hence the circle action is Hamiltonian for our sign
convention, and the autonomous Hamiltonian $k/2$ has time-one map $g_Y$.

The downstairs half-turn fixes exactly the four centers.  At each center
there are complex coordinates in which the circle has weights
$(\alpha,\beta)$ with $\alpha,\beta\in\{1,-1\}$.  Regardless of these
signs, the half-turn is
\[
  (z_1,z_2)\longmapsto(-z_1,-z_2).
\]
The local blow-up consists of pairs $(\ell,v)$, where $\ell\subset\C^2$ is
a complex line and $v\in\ell$.  The lifted half-turn is
\[
  (\ell,v)\longmapsto(\ell,-v).
\]
It fixes exactly the zero section, acts identically on its tangent, and
acts by $-1$ on its normal line.  Off the exceptional curves the blow-down
is an equivariant isomorphism, so no additional fixed points occur.  The
symplectic normal assertion follows also from the eigenspace decomposition
in Proposition~\ref{prop:involution}.

The homology formula~\eqref{eq:blowup-homology} adds one free generator in
degree two for each point blown up.  Since the homology of
$\CP^1\times\CP^1$ is free with ranks $1,2,1$ in degrees $0,2,4$,
respectively, it gives~\eqref{eq:Y-homology}.
\end{proof}

For later use, we state explicitly the product construction.  Its input
is an ordinary Morse function; its critical points will become genuine
Hamiltonian fixed points.

\begin{corollary}\label{cor:product-realization}
Let $(X,\omega_X)$ be a connected closed simply connected K\"ahler manifold,
let $(Y,\omega_Y)$ be the surface in Proposition~\ref{prop:surface}, and give
$M=X\times Y$ the product symplectic form
$\omega_X\oplus\omega_Y$.  Put $C=X\times\CP^1$.
For every smooth Morse function $h:C\to\R$, there is a smooth one-periodic
Hamiltonian $H$ on $M$ such that every fixed point of its time-one map is
nondegenerate and contractible for its Hamiltonian path, and
\[
  \#\Fix(\phi_H^1)=\#\Fix_0(\phi_H^1;H)
       =4\,\#\operatorname{Crit}(h).
\]
\end{corollary}

\begin{proof}
The Hamiltonian involution $g=\mathrm{id}_X\times g_Y$ fixes exactly the four
symplectic submanifolds $X\times E_j$, each diffeomorphic to $C$.  Transport
$h$ to each of them and apply Proposition~\ref{prop:involution}.  The product
$M$ is simply connected, so all the fixed-point loops are contractible.
More concretely, use the path generated by the pullback of $k/2$ in the
proof of Proposition~\ref{prop:surface}.  For a resulting fixed point
$(x,y)\in X\times E_j$, the circle part of the path lies in $\{x\}\times E_j$ and the
perturbation part is stationary.  The whole loop therefore contracts in
that projective line.
\end{proof}
